%-----------------------------------------------------------
%   Capítulo 7 - Resultados
%-----------------------------------------------------------
\chapter{Resultados}
\label{cap:resultados}

A placa ainda não foi fabricada nem o \textit{firmware} correu numa placa.
Este capítulo apresenta, por isso, os resultados da verificação do projeto,
que já existem, e os ensaios planeados para quando houver placa.

\section{Verificação do projeto}

A Tabela~\ref{tab:res-verificacao} junta as verificações que a CI corre em
cada alteração do repositório.

\begin{table}[htbp]
\centering\small\setstretch{1.05}
\renewcommand{\arraystretch}{1.15}
\caption{Verificações automáticas do projeto (26 de novembro de 2026).}
\label{tab:res-verificacao}
\begin{tabularx}{\textwidth}{@{}>{\raggedright\arraybackslash}X>{\raggedright\arraybackslash}p{0.42\textwidth}@{}}
\toprule
\rowcolor{LasetTint}
Verificação & Resultado \\
\midrule
ERC do esquema (KiCad 9) & 1 erro aceite (pinos VCAP unidos) \\
DRC da placa (KiCad 9) & 0 violações de cobre; 0 ligações em falta \\
Paridade entre esquema e placa & 0 diferenças \\
Testes dos módulos do \textit{firmware} no PC & 101 verificações, 0 falhas \\
Pinos do \textit{firmware} contra o esquema & 35 pinos, 0 diferenças \\
Compilação para NUCLEO e PCB & sem avisos \\
\bottomrule
\end{tabularx}
\end{table}

Algumas medidas da placa, retiradas do ficheiro do KiCad:

\begin{itemize}
  \item comando das portas dos MOSFET: 17 a 22~mm por porta, com o
        \textit{bootstrap} a 6 e 10~mm;
  \item condensadores do VCAP a 1,0 e 1,5~mm do pino; condensadores de
        desacoplamento dos circuitos integrados, na maioria, a menos de 5~mm;
  \item 452 vias (359 de 0,6/0,3~mm e 93 de 0,5/0,2~mm), nenhum par de furos
        de redes diferentes a menos de 0,5~mm;
  \item sinais na camada interior de 3,3~V: só sinais lentos (UART do GPS,
        CAN1\_TX, RS485\_TX, I2C\_OE e RBIAS).
\end{itemize}

O custo das peças da placa é de 118,99~€ por unidade, e o do projeto
completo, com a NUCLEO, o recetor GPS, o fabrico e os cabos, de 241,44~€.

\section{Ensaios planeados}

Os passos do primeiro arranque estão em \texttt{hardware/ARRANQUE.md}. A
Tabela~\ref{tab:res-ensaios} liga cada ensaio ao que verifica; os valores
medidos entram aqui depois dos ensaios.

\begin{table}[htbp]
\centering\small\setstretch{1.05}
\renewcommand{\arraystretch}{1.15}
\caption{Ensaios planeados para a placa.}
\label{tab:res-ensaios}
\begin{tabularx}{\textwidth}{@{}>{\raggedright\arraybackslash}p{0.24\textwidth}>{\raggedright\arraybackslash}X>{\raggedright\arraybackslash}p{0.17\textwidth}@{}}
\toprule
\rowcolor{LasetTint}
Ensaio & Como & Resultado \\
\midrule
Alimentação & 5, 12, 24 e 48~V à entrada; medir 5~V e 3,3~V, corrente em vazio & por medir \\
Estabilidade da fonte & degrau de carga de 0,2 para 1~A a 5 e 48~V; ondulação & por medir \\
Programação & SWD; LED a piscar a 2~Hz & por medir \\
RS232 e Ethernet & versão pela ferramenta do PC; \textit{ping} & por medir \\
CAN FD & tramas a 1 + 8~Mbit/s com um adaptador ou outra placa & por medir \\
RS485, I2C, PWM, PPS & ensaio de cada interface (Tabela em \texttt{ARRANQUE.md}) & por medir \\
Latência & GPIO de marcação e analisador lógico: média, p99 e máximo & por medir \\
\bottomrule
\end{tabularx}
\end{table}
